\chapter{Discussion}
\label{sec:discussion}

\section{\rqnolink{rq:comparison}: Comparing the three approaches}
\label{sec:rq1}

No approach wins the comparison outright:
a 139\gls{million}-parameter \gls{seqseq} model fine-tuned on a task-matched parallel corpus scores clearly and significantly higher than the untuned checkpoint on \gls{sari} and \gls{dsari} and roughly matches an existing simplifier three times its size on \gls{sari} at sentence level.
It is beaten by prompted instruction-tuned \glspl{llm} on \gls{sari} at sentence level and on \gls{lens} at document level.

The comparison therefore supports a claim about serving cost, and reinforces the methodological findings regarding \gls{ats} metric disagreement, rather than establishing a definitive ranking.
This is an issue analysed further in the cross-cutting findings (\cref{sec:cross-cutting}).

How the approaches rank depends on which systems are compared and which metric is used.
\Cref{tab:scoreboard} gives the four pairwise readings; they do not all agree in sign, and none supersedes the others.

\begin{table}[htbp]
	\centering
	\footnotesize
	\setlength{\tabcolsep}{3.5pt}
	\caption{Score Differences Between the Fine-Tuned Checkpoint and the Other Systems}
	\label{tab:scoreboard}
	\begin{tabular}{@{}P{3.1cm}P{4.25cm}P{4.8cm}@{}}
		\toprule
		\textbf{Comparison}                                 & \textbf{Sentence level}                              & \textbf{Document level}                  \\
		\midrule
		untuned \newline checkpoint                         & \(+16.42\)~\gls{sari}, \gls{ci}~\([+15.38, +17.50]\) & \(+21.09\)~\gls{dsari} (35.44 vs. 14.34) \\
		comparison \newline simplifier                      & \(-0.37\)~\gls{sari}, intervals overlapping          & not measured                             \\
		\gls{llm}, \newline \gls{ngram} metrics             & \(-8.15\) to \(-9.23\)~\gls{sari}                    & \(+12.54\)~\gls{dsari}                   \\
		\gls{llm}, learned human- \newline judgement metric & not measured                                         & \(-23.80\)~\gls{lens}                    \\
		\bottomrule
	\end{tabular}
	\note{
		Positive values mean that the fine-tuned checkpoint scores higher than the system in the row.
		Only the first row was tested for significance (paired bootstrap, \(\gls{pval} < 0.001\) for both margins).
		The \gls{llm} rows span the 7\gls{billion} and the 3\gls{billion} at sentence level and are the 7\gls{billion} at \(\gls{nsize} = 2{,}000\) at document level.
		\emph{Not measured}: the comparison simplifier has no document counterpart, and \gls{lens} was not computed at sentence level (\cref{sec:limitations}).
	}
\end{table}

\pagebreak

\paragraph{Against the untuned checkpoint, fine-tuning raises \gls{sari} and \gls{dsari} decisively.}
The sentence-level margin is robust to the one contamination worry the benchmark carries: eight ASSET test sentences also occur in the training data, and removing them moves it by \num{0.006}.
This comparison is the only one of the four that holds architecture and scale fixed, so it isolates the effect of fine-tuning; since the margin is significant at both granularities, Hypothesis~\cref{hyp:finetuning} holds.

\paragraph{Against the larger comparison simplifier, the \gls{sari} difference is within the noise floor, and the advantage is serving cost.}
The scores are \num{37.80} against \num{38.17}~\gls{sari}, a difference well inside the \({\approx}2\)-point noise floor estimated from two unseeded training runs (\cref{sec:sentence-results}).
No paired test between the two was run, but their intervals against the untuned checkpoint overlap almost completely.

The fine-tuned model reaches that score with 139\gls{million} parameters against 406\gls{million}, and on the one article profiled runs \num{2.4}~times faster per request and \num{2.8}~times faster per page (\cref{sec:runtime}).
On-device simplification faces the same constraint: \textcite[105]{romero-etal-2025-efficient} quantise their models, reducing their size by up to \qty{75}{\percent}, so that simplification runs entirely on the user's device.

Serving cost is also what links \cref{rq:comparison} to \cref{rq:deployment}, since the extension runs all of its models locally (\cref{sec:architecture}).
The comparison simplifier was fine-tuned on a pre-cleaned WikiLarge mirror \parencite[3]{cohen2026simplify}.
The fine-tuned model used the same mirror with \qty{1.8}{\percent} of its rows removed by the targeted filter (\cref{sec:corpus-preparation}), so the two checkpoints were trained on essentially the same data.
The comparison simplifier was scored without the \texttt{Simplify:} prefix its training used, but the score its authors report using the prefix is \num{0.21}~\gls{sari} lower than the one measured here, although with different scoring and decoding settings (\cref{sec:limitations}).

\paragraph{Against the prompted \glspl{llm}, the fine-tuned checkpoint loses on \gls{sari} and \gls{lens} and wins on \gls{dsari}.}
It loses at sentence level on \gls{sari}, and at document level on \gls{lens}, which was trained on sentence-level human judgements \parencite[16385--16386]{maddela-etal-2023-lens} and is applied here out of domain.
At document level it wins on \gls{dsari}, an \gls{ngram} metric \parencite[8002]{sun-etal-2021-document} computed here against a single reference.
Both document orderings hold across three sampling seeds, whose standard deviations of \num{0.071}~\gls{lens} and \num{0.039}~\gls{dsari} points are far below the gaps of \num{23.80} and \num{12.54} points.
This shows that neither ordering is a sampling artefact.
At document level, which system scores higher therefore depends on the metric.

\pagebreak

Hypothesis~\cref{hyp:llm-vs-finetuned} draws on the finding of \textcite[13291]{kew-etal-2023-bless} that their best prompted \glspl{llm} perform comparably with trained baselines.
The hypothesis, that a prompted \gls{llm} matches or exceeds the fine-tuned model, holds at sentence level: both prompted \glspl{llm} score above the fine-tuned checkpoint on \gls{sari}.
At document level the evidence is contradictory: the prompted 7\gls{billion} scores below the fine-tuned checkpoint on \gls{dsari} and above it on \gls{lens}, so the hypothesis remains open there.

\paragraph{The disagreement is a finding.}
Measurements rule out four artefacts as its cause: truncation, sample size, normalisation, and sampling seed.
The differential test against the reference implementation (\cref{sec:evaluation-protocol}) also rules out a defect in the \gls{dsari} reimplementation.
What remains is a disagreement with a plausible mechanism: \gls{dsari} rewards agreement with a single reference and penalises structural departure from it, whereas \gls{lens} is trained to reproduce human ratings of simplified sentences.

\textcite[148, 150--151]{alva-manchego-etal-2020-survey} report that \gls{bleu} and \gls{sari} correlate with human judgement only on some dimensions, and poorly once sentences are split.
Here the two metrics diverge far enough to reverse a system ranking.
One unmeasured factor could also affect the prompted \gls{llm}: in the document evaluation it receives the D-Wikipedia test sources as they are, lowercased and pre-tokenised, unlike the natural prose it receives in the extension.
If this input format changes its scores at all, it more plausibly lowers than raises them.

The disagreement cannot be settled here: \gls{lens} is applied outside the sentence-level data it was trained on (\cref{sec:limitations}), so it cannot act as the tie-breaker, and this project has no human ratings with which to decide between the two metrics.
The document-level system ordering under \cref{rq:comparison} is therefore metric-dependent and remains open; \cref{sec:future-work} specifies a study that could settle the disagreement.

\paragraph{Two qualifications limit the gain over the untuned checkpoint.}
First, \gls{bleu} favours the untuned checkpoint, which leaves \qty{93.9}{\percent} of sentences unchanged against the fine-tuned checkpoint's \qty{11.7}{\percent}.
This is expected, since \gls{bleu} does not compare the output with the input and so favours conservative, low-edit output (\cref{sec:background-evaluation}).
Since \gls{sari} ranks the untuned checkpoint last, the two metrics rank the conditions differently, and Hypothesis~\cref{hyp:metrics} holds.

Second, the aggregate scores do not reveal at least one systematic class of factual error.
The reduced-scope document checkpoint rewrote the census year 2010 as 2000 in 27 of 35 affected documents (\cref{sec:qualitative}), and this was found only by inspecting the outputs, not from any score.
The improvement over the untuned checkpoint is therefore significant on these metrics, which do not show whether the outputs are factually faithful.

\section{\rqnolink{rq:deployment}: Effective strategies and behavioural robustness}
\label{sec:rq2}

The answer to \cref{rq:deployment} rests on an exploratory site audit of twelve pages labelled by the author, and \cref{sec:limitations} sets out what an audit of this scale cannot show.

\paragraph{\Gls{dom} preservation}, the part of \cref{rq:deployment} that concerns replacement, holds.
Across 24 runs on twelve authentic pages, no tag, image source, or link target was lost, and reverting restored the text content and attribute sets on every run.
The raw \gls{html} is byte-identical on 18 of the 24 runs.
The six differences are cosmetic: once the extension writes an inline style, \texttt{jsdom} re-serialises the page's own inline styles, while the computed styles stay the same.

\paragraph{Content selection mostly withholds the right text, with two diagnosed false-positive triggers.}
The filters withhold \qty{74.7}{\percent} of sentence-cut candidates and \qty{90.9}{\percent} of document-cut candidates (\cref{sec:deployment-results}), and 48 of the 69 hand-labelled skips are correct.
The 15 wrongly withheld items in the labelled sample are ordinary prose that the code-like-text heuristic rejected: 12 because it contains a semicolon and 3 because it contains the word ``let''.

Removing the semicolon from the heuristic is not a simple fix: it would release all \num{102} items the semicolon matched, not only the 17 running-prose items among them.
The audit was the last study before data collection closed, so the heuristic was left unchanged; \cref{sec:future-work} gives the trade-off.

\paragraph{Most emptied links and the negation deletions trace back to the model's edits rather than to the \gls{dom} machinery.}
Of the \num{331} links left without text, \num{237} lost it because the model deleted the phrase they were attached to.
Deletion is a regular simplification operation, and write-back can only place a link on text that survives the rewrite (\cref{sec:content-selection}).
The remaining 94 are link text misplaced by the fragmentation defect, 93 of them on a single densely linked article.
The negation deletions are model edits that the change guard lets through by design (\cref{sec:deployment-results}).
Fixes for these failures therefore belong mainly in the guards and the model, with the fragmentation defect as the one cause on the extension side.

The latency measurements add one further observation.
On one Wikipedia article, for example, the sentence cut sent 79 small requests (median \num{124} characters) and the document cut 11 large ones (median \num{1262} characters) for almost the same text.

\pagebreak

The sentence cut finished the page sooner (\qty{15.5}{\second} against \qty{22.5}{\second}), because its requests overlap in the client's eight-slot queue, but its latency varied more: its \gls{p95} was \(2.6\times\) its median, while the document cut's was almost equal to its median.
In a separate run made while the machine was also running an evaluation job, the sentence cut's wall clock was \num{7.8} times and its \gls{p95} latency 19 times as long, while the document cut's stayed about the same.
For an extension that shares a laptop with the user's other programs, this stability may matter more than speed, although it rests on a single run; a comparison of throughput alone would not have shown it.

Several deployment defects left no visible trace and were found only on closer inspection, by checking the page's structure after simplification or by checking that item counts add up.
Some examples are: a folded section that took a list of links with it; a log that recorded what was sent rather than what the page showed; leaves that went uncounted because the document cut registered only one item per section; and prose withheld on a semicolon without any indication to the reader (\cref{sec:deployment-results,sec:granularity-findings}).
Two of these are absences rather than wrong outputs, and no benchmark would show any of them.
Two findings of the audit were also invisible to earlier checks: the withheld prose needed a skip reason that no interface reported, and the negation deletions needed a word-by-word comparison of input and output.
CheckList starts from the observation that held-out accuracy often overestimates real performance \parencite[4902]{ribeiro-etal-2020-beyond}.
Most of the failures found here lie in parts of the system that no benchmark evaluates, such as content selection and write-back, so Hypothesis~\cref{hyp:web} holds, and in a stronger form: the benchmarks cannot observe these failures at all.

\section{Cross-cutting findings}
\label{sec:cross-cutting}

Three observations plausibly transfer beyond the particular models used here.
Each is supported by at least two measurements in this project, though none was isolated experimentally.

\paragraph{Training and deployment should agree on the unit of text.}
A mismatch arose at both granularities.

At sentence level it happened: selecting text node by text node kept inline formatting but split sentences, so the sentence-trained models received \num{173} units for one article's prose instead of 79 sentences (\cref{sec:granularity-findings}).

At document level it was only a risk: sending one element at a time would have given the document-trained model single paragraphs, which the heading-delimited document cut avoids (\cref{sec:unit-of-simplification}).

\pagebreak

Neither corpus nor benchmark could have revealed either mismatch, because both lie in the interface between the model and its input, not in the model or the data.
Document-level corpora, metrics, and models address the modelling side of this problem \parencites[7999, 8002]{sun-etal-2021-document}[1, 3]{vasquezrodriguez2024simple}, but not the deployment side.

\paragraph{Benchmark and metric choice can change conclusions}, as two results show.
First, \gls{bleu} ranks the untuned checkpoint above the fine-tuned sentence checkpoint on ASSET, but below the fine-tuned document checkpoint on D-Wikipedia.
The untuned model is the same in both evaluations, but the fine-tuned checkpoints differ, so the reversal cannot be assigned to the benchmark alone; a plausible contributor is how much copying the source is worth against each benchmark's references \parencite[404, 410]{xu-etal-2016-optimizing}.
Second, \gls{dsari} and \gls{lens} rank the fine-tuned document checkpoint and the prompted 7\gls{billion} in opposite orders, consistently across seeds.

\paragraph{Inspecting the outputs finds what the scores do not show.}
The structured inspection of generations already on disk (\cref{sec:qualitative}) found both the census-year substitution and the disambiguation stubs paired with full articles.
From the reduced-scope to the full-scope checkpoint, the rate of stating a wrong census year falls from \qty{77}{\percent} to \qty{0.2}{\percent}, while the corpus-level \gls{dsari} rises by \num{3.82}.
That score cannot show whether its rise includes this correction, since the affected documents make up only \qty{7.2}{\percent} of the test split and the metric is not broken down by error class.
This is the case for the human evaluation proposed in \cref{sec:future-work}, made from this project's own data.

The deployment showed the same problem: \gls{dom} integrity and linguistic integrity are separate properties, and a check on one says nothing about the other.
While the system was sending sub-sentence fragments to the model, the robustness check only tested whether the page's \gls{dom} structure survived.
It therefore kept passing, although the text the model received was linguistically broken.
